% ===================================================================
%  TAFEL Nr. 13 — Das Hotel. --- Das Restaurant. --- Das Café.
%  Traduction 2026 d'apres texto/io, controlee sur texto/fr et
%  texto/en. Consigne : voir texto/de/10-tafel-01.tex.
%
%  Deux notes, marquees toutes deux « (*) », et deux appels : celui du
%  bloc 01-2 (la chambre) et celui du bloc 09-1 (le menu). L'ordre les
%  departage.
%
%  C'EST LA PAGE LA PLUS FRANCAISE DES SEIZE, ET ELLE L'EST EXPRES.
%
%  Hotel, Portier, Etage, Salon, Restaurant, Café, Menü, Karte,
%  Service, Couvert, Entrée, Dessert, Buffet, Terrasse, Balkon,
%  Serviette, Chef, Billard : le personnel et les lieux d'un grand
%  hotel allemand de 1926 portent des noms francais d'un bout a
%  l'autre. Ce n'est pas faute de mots : l'allemand a « Gasthaus »,
%  « Kellner », « Speisekarte », « Nachtisch », « Stockwerk ». Il les
%  a, et l'hotel ne les emploie pas.
%
%  VOILA UN CINQUIEME MECANISME, ET LA SERIE NE L'AVAIT PAS ENCORE
%  RENCONTRE. Les quatre premiers etaient :
%     le mot arrive avec la chose      (Gletscher, tableau 9) ;
%     le mot est decrete               (Turnen, tableaux 2 et 10) ;
%     la forme est calquee             (Bahnhof, tableau 12) ;
%     l'institution garde la langue du pays qui l'a fondee
%                                      (tableau 11).
%  Le cinquieme est different de tous : ON GARDE LE MOT ETRANGER
%  PARCE QU'ETRE ETRANGER EST JUSTEMENT CE QU'IL SERT A DIRE. Le
%  garcon d'un grand hotel s'appelle « Garçon » et non « Kellner » non
%  pas parce qu'on manque du second, mais parce que le premier coute
%  plus cher. LE MOT NE NOMME PAS LA CHOSE, IL NOMME LE PRIX DE LA
%  CHOSE.
%
%  ET LA PREUVE EST DANS LA MEME PAGE, A L'ALINEA 9. Le menu est
%  francais — Menü, Karte, Entrée, Dessert — mais CE QUI EST DANS
%  L'ASSIETTE est allemand : Krabben, Butter, Kutteln, Hammelkeule,
%  Spinat, Käse, Obst, Gebäck, Wein. La categorie est francaise et ses
%  membres sont allemands, exactement comme le « Möbel » du tableau 6.
%  Deux fois la meme figure, pour deux raisons opposees : au tableau 6
%  c'etait le droit, ici c'est le prestige. La figure a donc un nom,
%  et il faudra la chercher ailleurs.
%
%  UN MOT DE CETTE PAGE PERMET D'ESSAYER LE CRITERE DU TABLEAU 12. Le
%  pourboire de l'alinea 9 est en allemand Trinkgeld,
%  l'argent-a-boire. Le tableau 13 estonien avait trouve
%  « jootraha » — argent-a-boire, morpheme pour morpheme. Le francais
%  dit « pourboire », la meme image encore ; l'anglais dit « tip », et
%  ne la partage pas.
%    L'image est-elle evidente ? A demi : on comprend, mais on ne la
%    trouve pas seul, sinon l'anglais l'aurait. Le critere du tableau
%    12 laisse donc la question ouverte, et l'histoire la tranche :
%    l'estonien a pris son mot a l'allemand, comme il lui a pris le
%    tiers de son lexique, tandis que le « pourboire » francais est
%    atteste au XVIIe siecle, avant tout contact. UN CALQUE D'UN COTE,
%    UNE CONVERGENCE DE L'AUTRE, POUR LE MEME MOT.
% ===================================================================

%%K t13-noto-1 noto
\VUnotes{143.2mm}{%
\textit{\VUgras{(*) Zu den Einzelheiten des Schlafzimmers siehe Tafel Nr. 6.}}%
}

%%K t13-apar-1 apar
\VUsaut{3.0mm}
\VUtitre{15.0pt}{60}{TAFEL Nr. 13}
\VUsaut{1.6mm}
\VUpk{10.2pt}{40}{Das Hotel. --- Das Restaurant.}
\VUsaut{2.0mm}
\VUpk{10.2pt}{40}{Das Café.}
\VUsaut{3.4mm}

%%K t13-01-1 p
1. --- Ich bin gestern Abend in Royan angekommen und bin im \textit{Hotel Ozean}
abgestiegen, das mir ein Freund empfohlen hatte.

%%K t13-01-2 p
Man hat mir ein schönes \VUgras{Zimmer} \textsuperscript{(2)} im zweiten
\VUgras{Stock} \textsuperscript{(1)} gegeben, wo ich sehr gut geschlafen habe (*).

%%K t13-02-1 p
2. --- Als ich heute Morgen mein Zimmer verließ, wo mir das \VUgras{Zimmermädchen}
\textsuperscript{(3)} das Frühstück gebracht hatte, ging ich in das
\VUgras{Rauchzimmer} \textsuperscript{(5)}, das zugleich als \VUgras{Lesezimmer}
\textsuperscript{(4)} dient, um die \VUgras{Zeitungen} \textsuperscript{(6)} bei einer
\VUgras{Zigarette} \textsuperscript{(8)} durchzusehen. Als ich fertig gelesen hatte,
fuhr ich mit dem \VUgras{Fahrstuhl} \textsuperscript{(9)} hinunter und ging in der
Stadt spazieren.

%%K t13-03-1 p
3. --- Weil ich vergessen hatte, den \VUgras{Angestellten} \textsuperscript{(12)} des
Hotels zu fragen, wann an der \VUgras{Table d'hôte} \textsuperscript{(11)} gegessen
wird, kam ich früh zurück und nutzte die Gelegenheit, im \VUgras{Badezimmer}
\textsuperscript{(10)} des Hotels ein Bad zu nehmen. Dann ging ich zur
\VUgras{Kasse} \textsuperscript{(15)}, um einen Freund anzurufen. Aber das
\VUgras{Telefon} \textsuperscript{(13)} war besetzt; als die \VUgras{Kassiererin}
\textsuperscript{(14)}, die gerade eine \VUgras{Rechnung} \textsuperscript{(17)} in
das \VUgras{Fremdenbuch} \textsuperscript{(16)} eintrug, sah, dass ich nicht weiter
wusste, bat sie den \VUgras{Portier} \textsuperscript{(18)}, den \VUgras{Hausdiener}
\textsuperscript{(19)} mit seinem \VUgras{Fahrrad} \textsuperscript{(20)} für mich
losschicken zu lassen.

%%K t13-04-1 p
4. --- Ich dankte ihr und ging hinaus, um dem \VUgras{Dienstmann}
\textsuperscript{(24)} (dem Gelegenheitsarbeiter) ein Trinkgeld zu geben, der auf seinem
\VUgras{Handkarren} \textsuperscript{(25)} meine \VUgras{Hutschachtel}
\textsuperscript{(26)}, meinen \VUgras{Handkoffer} \textsuperscript{(27)} und meine
\VUgras{Reisetasche} \textsuperscript{(28)} vom Bahnhof gebracht hatte. Der
\VUgras{Briefträger} \textsuperscript{(21)}, der gerade kam, reichte mir einen
\VUgras{Brief} \textsuperscript{(22)}.

%%K t13-05-1 p
5. --- Ich blieb noch eine Weile auf der Türschwelle neben dem \VUgras{Briefkasten}
\textsuperscript{(23)} stehen und sah dem Verkehr auf der Straße zu. Der
\VUgras{Kutscher} \textsuperscript{(30)} des \VUgras{Omnibusses}
\textsuperscript{(29)} lud gerade den \VUgras{Koffer} \textsuperscript{(31)} eines
\VUgras{Gastes} \textsuperscript{(32)} auf seinen Wagen, von dem der
\VUgras{Hotelbesitzer} \textsuperscript{(33)} sich verabschiedete. Ein junger
\VUgras{Telegrafenbote} \textsuperscript{(35)} brachte ohne jede Eile ein
\VUgras{Telegramm} \textsuperscript{(34)} für einen der Gäste des Hotels; ein
\VUgras{Tourist} \textsuperscript{(36)} mit seinem \VUgras{Fotoapparat}
\textsuperscript{(38)} über der Schulter sah in seinem \VUgras{Reiseführer}
\textsuperscript{(37)} nach.

%%K t13-tit-1 sub
\VUsaut{4.0mm}
\VUpk{10.2pt}{40}{Essenszeit.}
\VUsaut{3.0mm}

%%K t13-06-1 p
6. --- Vor dem Gitter döste ein gemütlicher \VUgras{Lastträger}
\textsuperscript{(39)} zwischen seinen \VUgras{Traghaken} \textsuperscript{(40)} und
seinem \VUgras{Schuhputzkasten} \textsuperscript{(41)}, während eine junge
\VUgras{Wäscherin} \textsuperscript{(42)} mit einem \VUgras{Korb}
\textsuperscript{(43)} Wäsche über die feierliche Miene des \VUgras{Oberkellners}
\textsuperscript{(44)} lachte, der an der Tür stand.

%%K t13-07-1 p
7. --- Der Anblick des \VUgras{Kochs} \textsuperscript{(45)}, der aus seiner
\VUgras{Küche} \textsuperscript{(47)} im \VUgras{Souterrain} \textsuperscript{(48)}
gekommen war, um einen \VUgras{Küchenjungen} \textsuperscript{(46)} auf einen Gang zu
schicken, erinnerte mich daran, dass es bald Essenszeit war, und ich ging über den
Hof zum Restaurant, wo ein \VUgras{Autofahrer} \textsuperscript{(50)} gerade sein
\VUgras{Auto} \textsuperscript{(49)} abstellte, während ein \VUgras{Mechaniker}
\textsuperscript{(52)} nach den Anweisungen des \VUgras{Radfahrers}
\textsuperscript{(53)} ein \VUgras{Motordreirad} \textsuperscript{(51)} reparierte,
das gerade einen Unfall gehabt hatte.

%%K t13-08-1 p
8. --- Ich fragte einen jungen \VUgras{Pikkolo} \textsuperscript{(54)}, der
\VUgras{Gläser Bier} \textsuperscript{(55)} trug, ob die Table d'hôte unter der
Veranda sei, und als er ja sagte, setzte ich mich an einen der Tische und aß
ausgezeichnet.

%%K t13-noto-2 noto
\VUnotes{143.2mm}{%
\textit{\VUgras{(*) Mit Hilfe der Speisenliste im Wörterverzeichnis kann der Lehrer
das folgende Menü leicht abwandeln.}}%
}

%%K t13-tit-2 sub
\VUsaut{4.0mm}
\VUpk{10.2pt}{40}{Ein ausgezeichnetes Mittagessen.}
\VUsaut{3.0mm}

%%K t13-09-1 p
9. --- Der Kellner brachte mir die Speisekarte (*) und die Weinkarte. Ich wählte und
bestellte \VUgras{Krabben} mit \VUgras{Butter} als \VUgras{Vorspeise} und als
\VUgras{ersten Gang} \VUgras{Kutteln auf Caener Art}. \VUgras{Fisch} nahm ich
keinen, aber ich bat um eine Portion \VUgras{Hammelkeule} und als \VUgras{Gemüse}
\VUgras{Rahmspinat}. Weil mich dann keine der \VUgras{Süßspeisen} reizte, ließ ich
mir allerlei Nachtisch bringen: \VUgras{Käse}, \VUgras{Obst} und \VUgras{Gebäck}.
Dazu nahm ich einen ausgezeichneten Saint-Émilion-\VUgras{Wein}, und sehr zufrieden
mit meinem Essen stand ich vom Tisch auf, nachdem ich dem \VUgras{Kellner}
\textsuperscript{(57)} ein ansehnliches \VUgras{Trinkgeld} \textsuperscript{(59)}
gegeben hatte.

%%K t13-10-1 p
10. --- Als der \VUgras{Geschäftsführer} \textsuperscript{(60)} sah, dass ich gehen
wollte, schlug er mir vor, auf den Balkon zu treten und das großartige Panorama des
\VUgras{Ozeans} \textsuperscript{(62)} zu bewundern. In der Ferne konnte man kleine
Fischer\VUgras{boote} \textsuperscript{(63)}, \VUgras{Segelschiffe}
\textsuperscript{(64)} und einen riesigen \VUgras{Ozeandampfer}
\textsuperscript{(65)} erkennen, dessen schwarze Silhouette sich von den weißen
\VUgras{Wolken} \textsuperscript{(67)} am \VUgras{Horizont} \textsuperscript{(66)}
abhob. Ein leichter Wind bewegte die \VUgras{Fahne} \textsuperscript{(70)}, die über
dem \VUgras{Schild} \textsuperscript{(69)} der \VUgras{Halle}
\textsuperscript{(68)} steckte.

%%K t13-tit-3 sub
\VUsaut{3.0mm}
\VUpk{10.2pt}{40}{Vergnügungen.}
\VUsaut{3.0mm}

%%K t13-11-1 p
11. --- Der Anblick war so fesselnd, dass ich beschloss, am nächsten Tag auf die
Terrasse des Cafés hinaufzugehen und ein \VUgras{Opernglas} \textsuperscript{(71)}
oder ein \VUgras{Fernrohr} \textsuperscript{(72)} mitzunehmen.

%%K t13-12-1 p
12. --- Als ich vom Balkon kam, ging ich ins Café des Hotels, auf ein
\VUgras{Getränk} \textsuperscript{(73)} und eine Partie \VUgras{Billard}
\textsuperscript{(75)}. Ich ging quer durch den Gastraum, wo viele Gäste
\VUgras{Schach} \textsuperscript{(78)}, \VUgras{Dame} \textsuperscript{(79)},
\VUgras{Domino} \textsuperscript{(80)}, \VUgras{Puff} \textsuperscript{(81)} oder
\VUgras{Karten} \textsuperscript{(82)} spielten, und stieg geradewegs ins
\VUgras{Billardzimmer} \textsuperscript{(74)} hinauf.

%%K t13-13-1 p
13. --- Als die Partie zu Ende war, ging ich ins Erdgeschoss hinunter und bat den
Kellner um einen \VUgras{Umschlag} \textsuperscript{(84)}, etwas \VUgras{Papier}
\textsuperscript{(83)}, eine \VUgras{Briefmarke} \textsuperscript{(85)}, eine
\VUgras{Feder} \textsuperscript{(87)} und etwas \VUgras{Tinte} \textsuperscript{(86)},
um einen Brief zu schreiben.

%%K t13-13-2 p
Dann rief ich einen \VUgras{Kutscher} \textsuperscript{(92)}, dessen \VUgras{Wagen}
\textsuperscript{(88)} vor dem Café gehalten hatte. Ich fragte ihn nach seinem Preis
für die Stunde und für die Fahrt, und als wir uns einig waren, öffnete er mir den
\VUgras{Wagenschlag} \textsuperscript{(89)} und setzte mich hinein. Er stieg wieder
auf seinen \VUgras{Bock} \textsuperscript{(91)}, trieb die Pferde mit seiner
\VUgras{Peitsche} \textsuperscript{(93)} kräftig an, und los ging es zu einer
herrlichen Fahrt.
\VUsaut{6.0mm}
\VUfilet{22mm}
